\section{已实现功能}
本项目面向校园通勤中的乘车预约、司机执行任务和车辆调度需求，采用 Flutter 单工程构建普通用户端、司机端与后台管理端。现阶段已完成角色入口、核心数据模型和三端业务原型，能够在内置演示数据上运行主要操作流程。小组目前已实现的功能及责任人如表\ref{tab:functions}所示。

\begin{table}[H]
\centering\small
\caption{已实现功能列表}\label{tab:functions}
\begin{tabularx}{\linewidth}{|P{2.0cm}|Y|P{1.6cm}|}
\hline
\centering 类别 & \centering 功能 & \centering\arraybackslash 责任人\\\hline
公共入口 & 演示账号登录、入口切换、角色匹配与路由重定向 & 许盛凯\\\hline
普通用户端 & 浏览车次、查看余座、选择座位并创建预约 & 许盛凯\\\hline
普通用户端 & 查看与筛选个人预约、取消预约和信用扣分 & 许盛凯\\\hline
普通用户端 & 模拟支付、信用等级展示、本地通知中心 & 许盛凯\\\hline
司机端 & 查看派车任务、乘客名单和上车点 & 赵凯彭\\\hline
司机端 & 开始任务、到站完成、单个活动任务约束 & 赵凯彭\\\hline
司机端 & 每5秒生成模拟位置记录、展示位置共享状态 & 赵凯彭\\\hline
司机端 & 日、周、月维度切换及预置任务统计展示 & 赵凯彭\\\hline
后台管理端 & 车辆与司机信息维护、绑定校验和删除限制 & 许盛凯\\\hline
后台管理端 & 人工调度、规则辅助建议、人工确认生效 & 许盛凯\\\hline
后台管理端 & 模拟车辆位置、示例 ETA 与历史分析面板 & 许盛凯\\\hline
工程验证 & 预约、任务、调度等关键规则的测试与整理 & 许盛凯、赵凯彭\\\hline
\end{tabularx}
\end{table}

目前界面层通过 Riverpod 读取和更新业务状态，GoRouter 管理三类工作台的访问入口，Mock Repository 提供初始车次、任务、车辆和需求数据。模型、状态控制和页面展示已分离，后续可以在保留主要界面交互的基础上替换数据来源。

当前阶段成果以可运行原型为范围。支付操作更新本地支付状态，位置共享使用定时生成的坐标，后台辅助调度采用规则匹配；各端演示数据尚未通过统一服务端同步。数据持久化、真实支付、GPS 采集和实时通信是下一阶段需要继续设计与联调的内容。
